% Unit under test. Identification data comes from metadata.tex.

\FunctionalDescription{
  \TheUnit\ stores the heater temperature setpoint, checks it against the
  allowed range 0--100 and restores the last valid setpoint from NvM at
  start-up. Details: \DocRef{design}.
}

% {name}{direction}{type / range}{description}
\Interface{SetpointCtrl_Init()}{in}{void}{Initialises the unit, reads NvM}
\Interface{SetpointCtrl_Set(sp)}{in}{uint8, [0, 100]}{Sets a new setpoint, returns E_OK / E_NOT_OK}
\Interface{SetpointCtrl_GetState()}{out}{enum}{IDLE / DEGRADED}

% [component]{dependency}{replaced by}{behaviour in test}
% [component] is optional: functions of one component share a box in the diagrams.
\Dependency[NvM]{NvM_ReadBlock()}{Mock}{Returns a configurable value, records calls}
\Dependency[NvM]{NvM_WriteBlock()}{Mock}{Records written data}
\Dependency[Dem]{Dem_ReportError()}{Mock}{Records reported events}

% Who uses the unit: {component}{interfaces it calls}
% Callers + dependencies give the context diagram; dependencies + interfaces
% give the test harness diagram in the Test Environment chapter.
\Caller{HeaterApp}{SetpointCtrl_Init(), SetpointCtrl_Set()}
\Caller{DiagService}{SetpointCtrl_GetState()}

% Free text: any LaTeX, printed at the end of this chapter (Unit Under Test).
% Write \end{FreeText} on its own line.
\begin{FreeText}
\subsection{Public API (excerpt)}
From \texttt{include/setpoint\_ctrl.h}, copied to the unit's \texttt{code/} folder:
\InputCode{setpoint_ctrl.h}
\end{FreeText}

